They came out of the tree line an hour before dusk, and they did not charge.

Jon had spent four days building himself ready for the other thing. He had stood on the wall-walk above the King’s Gate at every grey dawn since the column came inside these walls and rehearsed it in his own head the way a man rehearses a wound he knows is coming — the horns, the first shapeless press against the outer trench, the moment when the fighting stopped being a plan and became simply whatever each man could manage in the six feet of ground he was standing on. He knew the shape of that. He had lived inside it twice already, at the Wall and at Winterfell, and he had come to understand that the terror of it was at least a terror with edges.

This had no edges at all.

They filled the tree line from the northern road to the river bend and then they stopped, ninety yards past the furthest reach of any bow on that wall, and they stood.

“Why aren’t they moving,” said the boy beside him — a Northern levy, sixteen at the outside, holding a dragonglass-headed spear with both hands the way a man holds something he has been told will save him and does not yet believe. “Ser. My lord. Why aren’t they moving?”

“I don’t know,” Jon said, because the boy deserved the truth more than he deserved comfort, and because a lie would have been found out inside the hour anyway.

\sceneBreak

Word went along the line the way it always did, faster than any runner. The dead had come. The dead had stopped. Nobody could say why.

Tyrion arrived on the wall-walk at full dark, out of breath from the stairs, and stood a long while looking out at the torchlit ground without saying anything at all, which from Tyrion Lannister was its own kind of report.

“How many,” he said finally.

“More than we counted at Winterfell. Grey Worm’s people have been working the arithmetic since they stopped.” Jon kept his eyes on the field. “It doesn’t matter much past a certain number, my lord. Past a certain number it’s simply more than we have.”

“It matters to the ledgers.” Tyrion’s voice had gone flat. “It always matters to the ledgers, which is the great useless curse of my particular skills.” He gripped the stone. “Jon. Why have they stopped?”

“Sansa says he’s letting us stand at arms all night. That the not-attacking is the attack.” Jon had thought about it for the better part of six hours by then and had not managed to improve on his sister’s answer. “Every man on this wall has marched six hundred miles and slept four nights in a row on frozen ground. Now they’ll stand to arms until dawn, because no one on this wall will dare stand down while that’s out there. Come morning we’ll be a day worse than we were at dusk and he’ll have spent nothing at all to do it to us.”

“That’s not a wight’s thinking. That’s a commander’s thinking.”

“Aye,” Jon said. “That’s what’s been sitting in my chest since the hour they stopped.”

Tyrion was quiet. Out on the field, nothing moved that either of them could see, though the torchlight along the trench-line reached far enough to show the front of the mass — a stillness so complete it made a man’s eyes invent motion in it, so that Jon had to keep deliberately looking away and back to be certain of what he was seeing.

“At Winterfell he threw them at us,” Tyrion said slowly, working it aloud the way he did. “Wave after wave, no economy at all, no regard for how many he spent because spending them cost him nothing. That was the one comfort in the whole ruinous business — that he was enormous and stupid, and that a stupid enormous thing can sometimes be out-thought by a clever small one.”

“He’s not stupid. I think he’s never been stupid.” Jon found the words were coming out of him more steadily than he’d expected, as though some part of him had already made peace with the conclusion and had only been waiting for someone to say it out loud to. “I think at Winterfell he didn’t need to be careful. He had a castle full of tired people and eight times their number, and he’d have taken it whatever he did. Here he’s looking at trenches he’s never seen before, and pipes running into them out of cisterns, and scorpion mounts on a wall that weren’t there a month ago.” He nodded out at the still, patient dark. “He’s not waiting to frighten us, my lord. Or not only that. He’s \textit{reading} it. He wants to know what we’ve built before he walks into it.”

“Gods,” Tyrion said quietly, after a moment. “He learns.”

“Aye.”

“That’s considerably worse than more of them.”

“I know it.”

\sceneBreak

He walked the lines all night. There was nothing else useful left to do with himself, and he had learned at the Wall that a commander who stood still on a night like this became a statue men looked at instead of a man they could ask things of.

The Unsullied had their sections in perfect order, spears grounded, ranks dressed, nobody speaking. Grey Worm walked his own line the same way Jon walked his, and the two of them passed one another three times in the dark that night and exchanged perhaps twenty words between them all told, and Jon found the meeting steadied him more than a long conversation would have.

The Northern levies were worse. They were the ones who had marched out of Winterfell and they knew what was standing out there in a way the Unsullied and the Golden Company only knew secondhand. Jon sat with three of them near midnight and let them ask him things and answered honestly where he could, and lied exactly once, when a man from Long Lake asked whether it was true that a person who died out on that field would be up and walking by morning, and Jon said the dead needed time and the Night King’s attention both, which was not true and was the kindest thing available to him.

Roland Storm’s companies held the eastern stretch, and Jon stopped there too, because there was still a strangeness in it that he did not want to leave unaddressed — that these men had been holding a gate shut against him a week past and were now standing in his trench-line.

“They’ll fight,” Roland said, without being asked. “You’ll want to know that, and I’d not blame you for doubting it. They’ll fight, my lord, because there’s nowhere left on this continent to run to that’s better than behind these walls. That’s not loyalty. I’ll not sell you loyalty I can’t deliver. But it’s the next best thing to it, and in my experience it holds nearly as well while the fighting’s actually happening.”

“That’s more honest than most men would give me.”

“You’ll get a great deal of honesty out of me tonight, my lord. It’s cheap and I’ve nothing else spare.” Roland looked out at the field. “Can I ask you something plainly, since we’re both awake and neither of us is going to sleep?”

“Ask.”

“Is there any point at which a thing like that stops? Not stops being beaten — stops. Wants something, gets it, goes home.” Roland’s face in the torchlight was that of a man asking a professional question rather than a frightened one. “I’ve fought men my whole life. Every man I ever fought wanted something, and once he had it or gave up hope of it, he went home. I don’t know how to fight a thing that hasn’t got a home to go to.”

Jon thought of the heart tree at Winterfell and the dark sap on the carved face and Bran’s flat, careful voice saying \textit{he’s remembering something, or trying to forget it,} and found he had nothing he could give this man that would help him hold a trench.

“I don’t know,” he said. “My brother thinks it wants something. He can’t tell me what, and I’m not certain he’d be right if he told me.”

“Well.” Roland nodded slowly. “That’s honest too, at least.”

\sceneBreak

It was the hour past midnight when he saw them properly.

He had been avoiding it, if he was honest with himself, and had told himself for six hours that avoiding it was the same as having no cause to look. The torch-line along the outer trench threw light perhaps forty yards past the stakes, and the front of the mass stood at ninety, which meant the whole night had been a wall of shapes with nothing in it a man could resolve. Then somewhere along the middle stretch a fire-arrow went up — some levy’s nerve going, or an officer wanting the range confirmed — and it arced out and dropped and burned on the frozen ground well short, and for perhaps thirty seconds it lit the front rank from below.

Jon saw a man he had spoken to.

He did not get the name. That was somehow the worst of it, standing there afterward with his hands on the stone: he did not get the name, only the face, and the certainty that the face had been in front of him at some point in the last five weeks and had said something to him and he had answered it.

Then the arrow burned out.

He stood there a while. Below him on the wall-walk somebody was being quietly sick, and further along a sergeant was telling his men in a hard flat voice to keep their eyes front, which was good sergeanting and entirely useless.

The two hundred came into his mind then, and would not leave it.

He had gone to them himself. That had been the whole of what he’d insisted on, back at the narrows on the eleventh day, when Davos had offered to draw the lots and Sansa had offered to have the order written and sealed so it came from the crown rather than the man — he had said no to all of it and gone down the line himself and asked them to their faces, one after another, two hundred and some, because he had decided that a man who sends people to die owes them at minimum the discomfort of asking in person. More than two hundred had stepped forward. He had turned some away. He remembered turning away a grey-bearded man from Ramsgate who had argued with him about it, and he remembered thinking at the time that it was a mercy, and he understood now, standing on this wall, that the man he had turned away had marched on with the column and might have died anyway in some other ditch further south with nobody’s hand on his shoulder at all.

They were out there. Every one of them. He had asked them to hold a narrows and they had held it and their bodies had stayed where they fell because the column could not go back for them, and now they were standing in the front rank ninety yards from this wall, waiting to be told to come and kill the people they had died for.

\textit{That’s why they’re in front,} he thought.

He turned that over and could not make it dislodge. The dead had marched six hundred miles at the Night King’s pace and could have been arranged in any order at all, and the ones standing where the torchlight could almost reach them were the recent ones, the retreat’s own, the faces this particular army would recognize. It could be nothing. It could be simply that the newest dead walked at the back of the column and had arrived last and stood where they stopped.

He did not believe that.

\sceneBreak

Brienne found him near the third hour, and neither of them said anything for a long while, which he had come to appreciate about her.

“Ser Brienne.”

“My lord.” She stood at the stone beside him and looked out. Her jaw was set in the particular way he’d learned meant she was holding something at arm’s length by main force.

“Tormund’s out there,” Jon said. “Isn’t he.”

“I’ve not looked.” She said it evenly. “I’ve made a decision not to look, my lord, and I’d ask you not to help me break it, because I’ve work to do at dawn that wants my hands steady and I know myself well enough to know what looking would cost me.”

“Aye.” He was quiet a moment. “That’s better sense than I’ve managed tonight.”

“It isn’t sense. It’s cowardice with a job to do.” She almost smiled, and it did not reach anything. “He’d have laughed at me for it. He’d have said something unbearable about it and made me furious and then I’d have been fine for an hour.”

They stood together until the sky over the Blackwater began, very slightly, to go from black to the colour of dirty water.

\sceneBreak

Dawn came, and nothing came with it.

The light found them all still standing — a wall lined with men who had been at arms for thirteen hours, hands cramped around spears, eyes raw, the particular grey-faced hollowness that Jon knew from the Wall meant an army had already begun paying for a battle that had not yet started. Down in the trench-line the fires had burned low and nobody had been spared to feed them. Somewhere behind him in the city the bells rang the hour and half the wall flinched at it.

And ninety yards out, in full light now, with nothing left to hide behind, the dead stood exactly where they had stood at dusk and did not move.

Jon could see them clearly now. He made himself do it, all along the front, left to right, because it seemed to him that a man who had asked two hundred people to die for him did not get to look away from the result on the grounds that looking was hard. He found Robett Glover near the middle, which he had half expected and which was going to have to be kept from Rowan somehow, and he did not find Theon, and he did not find Podrick, and he understood that absence meant nothing except that he could not see everything.

He did not find Lyanna Mormont, and that one he was certain of, because they had burned her. They had burned every body the column had kept hold of, every night, however tired the men were and however much wood it cost them, because they had all learned that lesson early and none of them had ever needed teaching it twice. There was that, at least. Somewhere in all of it there was that.

“My lord.” A runner, breathing hard, at his elbow. “Lady Sansa’s asking for you. She says it’s urgent and she says it’s not the field.”

“Not the field.”

“She said to tell you that plainly, my lord. Not the field. Inside the walls.”

Jon looked out one last time at two hundred people he had asked a question of, standing patient in the morning light with their king’s business still to do, and turned away from the wall, and went to find out what his sister had learned that could possibly be more urgent than what was standing outside.